\documentclass[11pt]{article}
\usepackage[margin=1in]{geometry}
\usepackage{amsmath,amssymb,amsthm}
\usepackage{graphicx}
\usepackage{booktabs}
\usepackage{enumitem}
\usepackage{tabularx}
\usepackage{microtype}
\usepackage[hypertexnames=false]{hyperref}
\usepackage{algorithm}
\usepackage{algpseudocode}
\usepackage[nameinlink,noabbrev]{cleveref}

\newtheorem{theorem}{Theorem}
\newtheorem{lemma}{Lemma}
\newtheorem{proposition}{Proposition}

\title{Notarized Anonymity Certificates (Addendum):\\Extended Figures, Log Hardening Notes, and Templates}
\author{}
\date{April 27, 2026 (archive v764; support-carrier validation)}

\begin{document}
\maketitle

\begin{abstract}
This companion note collects material intentionally excluded from the brief ``Notarized Anonymity Certificates'' paper (Eval~1):
extra illustrative figures for grinding patterns, additional operational notes on log equivocation hardening (witnessing, gossip, cross-logging),
and short implementation-oriented templates. It intentionally stops at the supplementary evidence and policy-template layer:
it does not restate the core spend rule, retry-prefix verifier bundle, runtime conformance receipts, or release-lineage packaging.
All content is pure \LaTeX{} and remains source-only; supporting artifacts are available on request.\cite{series:artifactpolicy}
\end{abstract}

\paragraph{Scope.}
Addendum to Eval~1.\cite{series:evalnotary}

\paragraph{Imports.}
We import the core log-and-spend release discipline and manifest fields from Eval~1.\cite{series:evalnotary}
Transparency primitives (beacon/VRF binding, inclusion/consistency proofs) are owned by Synthesis~9.\cite{series:transparencyprimitives}

\paragraph{Exports.}
A compact supplementary packet: omitted-figure placeholders, log-hardening notes, a minimal checkpoint/proof sketch, and a minimal witness-policy template.

\paragraph{Earliest sufficient stop.}
Once a verifier can cite the supplementary packet as explanatory support for an Eval~1 claim, this addendum stops.
The core attempt-log/spend/verifier object remains in Eval~1; runtime-plan conformance receipts remain in Operational~B; and public lineage / release packaging remain in ABOM/OINL plus Release~A.\cite{series:evalnotary,series:opB,series:abomoinl,series:releaseA}

\paragraph{Later terse-paper citation posture.}
For later terse papers under space pressure, default to Eval~1 for the public retry-discipline / notarization claim and cite this addendum only when the exact supplementary packet exported here --- omitted-figure placeholders, log-hardening notes, checkpoint/proof sketch, or witness-policy template --- is itself live.
If the live question is only whether the evaluation used a non-gameable retry / abort discipline, stop at Eval~1.
If the live question is instead ``where is the explanatory support for that claim?'', widen to this addendum.
If the question becomes full audit, runtime-plan conformance, or public release-lineage comparison, leave this addendum and reopen the neighboring owners.
The maintained worked example and paired cutover rule are the general model for that stop / widen / reopen discipline.\cite{series:workedexample,series:pairedterminalcutoverrules}


\section{Exported object and intended use}
The exported object is deliberately small: a referee-readable supplementary packet that helps a reader understand how the main Eval~1 discipline behaves under grinding, witness equivocation, and practical checkpoint policies.
A public release can cite this packet for intuition, hardening notes, and lightweight implementation templates without promoting the addendum into a second theorem, notarization, or release-governance paper.

\paragraph{Stable handoff to neighboring layers.}
This addendum is easiest to audit when the stopping rule is stated directly rather than inferred from surrounding series context.
The handoff is therefore deliberately short and explicit.

\begin{center}
\begin{tabularx}{0.97\linewidth}{@{}>{\raggedright\arraybackslash}p{0.31\linewidth}>{\raggedright\arraybackslash}p{0.34\linewidth}>{\raggedright\arraybackslash}X@{}}
\toprule
\textbf{Referee question} & \textbf{Earliest sufficient stop} & \textbf{Owner} \\
\midrule
What extra support explains the retry-discipline claim? & This addendum's supplementary packet: omitted-figure placeholders, log-hardening notes, checkpoint/proof sketch, and witness-policy template. & Eval~1A \\
Was the evaluation itself run under a non-gameable retry / abort discipline? & Spend schedule, gap-free attempt prefix, and verifier report. & Eval~1 \\
Did the deployed traces still match the pinned runtime plan? & Plan certificate, PTL inclusion proof, and epoch conformance receipt. & Operational~B \\
How should this object be compared against earlier public releases? & Lineaged release receipt / release-ledger entry. & ABOM/OINL plus Release~A \\
\bottomrule
\end{tabularx}
\end{center}


That table is intentionally narrow: the addendum exports explanatory support for the notarization layer, not a competing certificate, runtime receipt, or publication object.

\paragraph{A minimal public notarization-support card.}
When a later short paper only needs to answer \emph{which explanatory support packet backs the published retry-discipline / anti-grinding claim?}, it should stop at the small card below plus the tiny witness-policy drill that follows.

\begin{table}[t]
\centering
\begin{tabularx}{0.97\linewidth}{@{}>{\raggedright\arraybackslash}p{0.31\linewidth}>{\raggedright\arraybackslash}X@{}}
\toprule
\textbf{Field} & \textbf{Declared support packet} \\
\midrule
Support packet id & \texttt{eval1a:supplementary-packet-v1}; the compact identifier for the explanatory packet itself. \\
Grinding figures & \texttt{release\_stop\_when\_pass}, \texttt{hyperparameter\_sweep\_grinding}, and \texttt{abort\_after\_beacon\_grinding}; the three omitted-figure placeholders that anchor the anti-grinding explanation. \\
Checkpoint hardening note & Signed checkpoints, append-only consistency, and witness / gossip requirements for split-view resistance. \\
Cross-log note & Independent-log submission is recorded as an extra audit surface rather than as a replacement for the main retry log. \\
Witness policy template & A local template only: \texttt{(W,q)=(7,4)} is kept here as a coarse sanity-check model, but the maintained worked bundle anchors the support packet through the concrete log/checkpoint/support fields below rather than by shipping that toy tuple as the public packet itself. \\
Main receipt hook & Retry / abort attempts count as spend and remain owned by Eval~1's attempt-prefix receipt, not by this support packet. \\
Release hook & Public lineage and package comparison remain owned by ABOM/OINL plus Release~A. \\
\bottomrule
\end{tabularx}
\caption{A minimal public notarization-support card. This is the smallest public answer later terse papers can quote directly when the live issue is the explanatory support packet behind Eval~1's retry-discipline claim rather than the main notarized receipt itself.}
\label{tab:eval1a-card}
\end{table}

\paragraph{Worked-bundle checkpoint-anchor drill.}
For the maintained worked bundle, the support-facing clause is no longer illustrative.
The addendum can now point to the concrete log/checkpoint/support route already carried by
\nolinkurl{example_uvi.json},
\nolinkurl{example_support_bundle_map.json}, and
\nolinkurl{example_replay_plans.json} in the worked-example artifacts.
That route names
\texttt{log\_id=example-transparency-log},
checkpoint \nolinkurl{tree_size=42 root=EXAMPLE_ONLY},
proof pointer \texttt{consistency://example-transparency-log/41-42},
receipt digest \texttt{095e828fd09ef186757f5dd49e84769603a2b59ea77c4a8bfb0edb720bc64cc1},
replay hook \texttt{eval-fallback-cct-v1} / \texttt{sha256:e9168dbec2cabd9115a7f3ea9aebe0a75ff910d2295e450a340716cc7e541193},
and support bundle \texttt{bundle://worked-example/fallback-cct} with support pointer \texttt{support://fallback/notarized-attempt-log.example}.
The maintained cut for this support carrier is \nolinkurl{worked-example-draft-199} / artifact note version \nolinkurl{1.99}; the same replay hook is pinned by plan spec \nolinkurl{sha256:e9168dbec2cabd9115a7f3ea9aebe0a75ff910d2295e450a340716cc7e541193}.
So the addendum's narrow hardening sentence is now: this support packet anchors the same maintained receipt at one concrete transparency checkpoint and one concrete on-request notarized-attempt-log pointer, while the retry / abort spend rule and verifier verdict still live in Eval~1.
The local witness-policy template remains useful only as a fail-closed explanatory sidecar; it is no longer the smallest honest shipped clause for the worked route.

\paragraph{Paired route back to the main notarized receipt.}
Suppose a later short note must say two things at once: first, that a headline anonymity number came from a logged retry-safe evaluation bundle; second, that the bundle's transparency anchor was supported by a concrete checkpoint-hardening packet.
The non-redundant route is to cite Eval~1 for the public replay object --- plan anchor, logged attempt prefix, many-testing rule, verifier threshold line, and transparency anchor --- and to cite Table~\ref{tab:eval1a-card} here only for the supplementary checkpoint-anchor / support-pointer sentence that explains why the log anchor is a credible one.
For example, a later terse note may honestly pair ``attempt $2$ was the first passing logged attempt under the disclosed retry budget'' with ``the same receipt was anchored at \texttt{example-transparency-log} / \texttt{tree\_size=42 root=EXAMPLE\_ONLY} and the worked support packet carried \texttt{support://fallback/notarized-attempt-log.example} for \texttt{eval-fallback-cct-v1}.''
But it may not replace the logged-prefix / verifier sentence with that checkpoint/support sentence, and it may not pretend that the support packet by itself settles runtime conformance or release lineage.

\paragraph{Evaluation-series four-stop route (support stop only).}
Within the local evaluation family this addendum is only the optional second stop: Eval~1 owns the replayable retry-safe evaluation object, this addendum owns the supplementary checkpoint / witness hardening packet, Eval~2 owns any later audit/live transfer sentence for a detectable monitor, and Eval~3 owns any later cadence-facing horizon or alert-tax conversion.
So a terse Anonymous-DHT note may widen from Eval~1 into this addendum when the support tuple itself matters, but it must keep moving onward to Eval~2 or Eval~3 when the live sentence has become a detectability or calibration sentence rather than a hardening sentence.

\paragraph{Same support packet / refreshed wrapper cutover.}
Suppose a successor release keeps the same concrete checkpoint-anchor/support tuple
\[
(\texttt{log\_id},\texttt{checkpoint},\texttt{proof\_pointer},\texttt{bundle\_id},\texttt{support\_pointer}).
\]
In the maintained worked route that means keeping
\texttt{example-transparency-log},
\nolinkurl{tree_size=42 root=EXAMPLE_ONLY},
\texttt{consistency://example-transparency-log/41-42},
\texttt{bundle://worked-example/fallback-cct},
and \texttt{support://fallback/notarized-attempt-log.example}
fixed while only the outer release wrapper, successor notice, or public carrier around the Eval~1 claim is refreshed.
Then a later terse paper may keep Table~\ref{tab:eval1a-card} unchanged and reopen only the release-facing owner chain.
If the live change is instead a new log id, checkpoint/proof pointer, support bundle, support pointer, or witness-policy / cross-log sidecar, the old support card is stale and the paper must widen back to this addendum (or to Eval~1 if the retry-prefix / spend rule changed too).

\section{Extra figures, hardening notes, and templates}

\begin{figure}[t]
\centering
\fbox{\begin{minipage}{0.92\linewidth}\centering\small Figure omitted in source-only release. Requested file: \texttt{figures/release\_stop\_when\_pass.png}.\end{minipage}}
\caption{Naive stop-when-pass release gates inflate false safety approvals (unsafe releases).}
\label{fig:stop_when_pass}
\end{figure}

Figure~\ref{fig:sweep_grinding} shows the same effect for metric/slice/seed sweeping: the naive false-approval rate approaches $1$ as the sweep library grows.

\begin{figure}[t]
\centering
\fbox{\begin{minipage}{0.92\linewidth}\centering\small Figure omitted in source-only release. Requested file: \texttt{figures/hyperparameter\_sweep\_grinding.png}.\end{minipage}}
\caption{Metric/slice/seed sweeping inflates false approvals unless budgeted.}
\label{fig:sweep_grinding}
\end{figure}

Figure~\ref{fig:abort} illustrates grinding in a simple model and how a log-enforced protocol blocks it.

\begin{figure}[t]
\centering
\fbox{\begin{minipage}{0.92\linewidth}\centering\small Figure omitted in source-only release. Requested file: \texttt{figures/abort\_after\_beacon\_grinding.png}.\end{minipage}}
\caption{Abort-after-beacon grinding can bias outcomes unless aborted attempts are logged and counted as spend.}
\label{fig:abort}
\end{figure}

Such statements can also coexist with SLSA provenance attestations when publishing software artifacts \cite{slsa_provenance_v1}.

Append-only Merkle logs are vulnerable to \emph{equivocation} (a \emph{split-world} attack): a malicious operator can present different, internally consistent log views to different verifiers.
Operational ecosystems mitigate this by (i) \emph{gossip} that compares signed tree heads/checkpoints across clients, and/or
(ii) \emph{witnessing} in which a diverse set of third parties cosign checkpoints so that a verifier can reject un-witnessed forks.

One deployable pattern is:
\begin{enumerate}
\item The log periodically emits a signed checkpoint (tree head + size).
\item Witnesses fetch checkpoints, verify append-only consistency, and cosign the latest checkpoint.
\item Verifiers require a threshold of witness signatures (or verify consistency via gossip/monitors) before accepting a checkpoint as authoritative.
\end{enumerate}
Witness cosigning can be implemented efficiently at scale (e.g., collective cosigning) \cite{syta2016cosi}, and CT-gossip variants have been studied for split-view detection \cite{dahlberg2018ctgossip}.

\paragraph{Cross-logging as an independent audit surface.}
An increasingly common hardening pattern is \emph{cross-logging}: periodically submit the log's signed checkpoint into an \emph{independent} log (or log pool), so that equivocation requires compromising \emph{multiple} operators and/or their monitoring infrastructure.
This is discussed broadly in systems work on log-based transparency \cite{hicks2023sok_transparency} and appears as a concrete design point in CT-adjacent proposals (e.g., pooled witnesses and log collaboration) \cite{dirksen2021logpicker}.
Figure~\ref{fig:crosslog} is a toy model showing how even a small cross-logging rate can sharply reduce the probability a split-view persists undetected across many checkpoints.

Figure~\ref{fig:split_view} is a simple simulation showing that even small witness sampling rates can make split-view attacks brittle when witnesses are diverse.

\begin{figure}[t]
\centering
\fbox{\begin{minipage}{0.92\linewidth}\centering\small Figure omitted in source-only release. Requested file: \texttt{figures/split\_view\_witness\_detection.png}.\end{minipage}}
\caption{Toy split-view model: detection probability rises quickly with witness diversity and modest sampling.}
\label{fig:split_view}
\end{figure}

\begin{figure}[t]
\centering
\fbox{\begin{minipage}{0.92\linewidth}\centering\small Figure omitted in source-only release. Requested file: \texttt{figures/cross\_logging\_robustness.png}.\end{minipage}}
\caption{Toy robustness comparison: witnesses-only vs witnesses plus cross-logging to an independent log/monitor.
Cross-logging forces an attacker to evade multiple audit surfaces, making long-lived split views brittle.}
\label{fig:crosslog}
\end{figure}


\paragraph{Cross-logging helps under partial witness compromise.}
Witnessing assumes diversity; in practice, some fraction of witnesses or monitors may be compromised or coerced.
A key benefit of cross-logging is that it introduces an \emph{independent} audit surface: even if some witnesses are silent,
a fork must evade both remaining honest witnesses \emph{and} the external log/monitoring path.
Figure~\ref{fig:crosslog_comp} illustrates this in a toy persistence model.

\begin{figure}[t]
\centering
\fbox{\begin{minipage}{0.92\linewidth}\centering\small Figure omitted in source-only release. Requested file: \texttt{figures/crosslog\_witness\_compromise.png}.\end{minipage}}
\caption{Toy split-view persistence with partial witness compromise.
Dashed curves add occasional cross-logging to an independent audit surface; detection becomes markedly more likely even when a large fraction of witnesses are compromised.}
\label{fig:crosslog_comp}

\end{figure}



\paragraph{A practical witness policy template.}
A referee will reasonably ask: \emph{how many witnesses, and what quorum?}
While the right answer depends on your operational assumptions, a useful \emph{sanity check} is the following conservative model:
if clients accept a checkpoint only when it carries at least $q$ cosignatures from a fixed set of $W$ recognized witnesses,
then sustaining \emph{two inconsistent} checkpoints that both validate requires (in a worst case) at least $2q$ compromised witnesses.
Thus a coarse feasibility bound is $\Pr[\mathrm{Bin}(W,p)\ge 2q]$ when each witness is compromised independently with probability $p$.
Figure~\ref{fig:witness_quorum} plots this bound for a few $(W,q)$ choices.
Notably, choosing $2q>W$ makes equivocation via witness compromise alone infeasible in this model, while still allowing high availability when witness downtime is modest.
In practice, witness protocols and public witness directories make it straightforward to configure a diverse witness set \cite{c2sp_tlog_witness}.

\begin{figure}[t]
\centering
\fbox{\begin{minipage}{0.92\linewidth}\centering\small Figure omitted in source-only release. Requested file: \texttt{figures/witness\_policy\_quorum\_tradeoff.png}.\end{minipage}}
\caption{Conservative split-view feasibility bound under witness compromise: $\Pr[\mathrm{Bin}(W,p)\ge 2q]$ for several witness policies $(W,q)$.}
\label{fig:witness_quorum}
\end{figure}

\section{Experiments (extended version)}
We include three evidence types.

\paragraph{Spending schedules are practical.}
Spending is not only safe but tunable. Figure~\ref{fig:tradeoff} compares a few standard summable schedules
(a $1/t^2$ schedule and geometric schedules) against a naive stop-when-pass gate.
Different schedules trade off speed (expected attempts until approval for a safe build) and conservatism,
while remaining within the same total false-approval budget $\delta_{\mathrm{tot}}$.

\begin{figure}[t]
\centering
\fbox{\begin{minipage}{0.92\linewidth}\centering\small Figure omitted in source-only release. Requested file: \texttt{figures/spending\_schedule\_tradeoff.png}.\end{minipage}}
\caption{Speed--safety tradeoff for common summable spending schedules (stylized).}
\label{fig:tradeoff}
\end{figure}

\paragraph{Tile-backed logs make monitoring affordable.}
A common objection to transparency-based release discipline is that independent parties will not \emph{actually} monitor the log at scale.
Modern deployments address this with tile-backed log APIs and compact proofs, where monitors fetch fixed-size tiles (amortizing repeated requests) rather than repeatedly downloading overlapping proofs. \cite{c2sp_tlog_tiles,static_ct_api_spec,tesseract_blog_2025,cloudflare_azul_ct_2025}
Figure~\ref{fig:tile_cost} is a stylized bandwidth model: as the number of audited entries grows, tile-backed auditing amortizes rapidly compared to per-entry inclusion proofs.

\begin{figure}[t]
\centering
\fbox{\begin{minipage}{0.92\linewidth}\centering\small Figure omitted in source-only release. Requested file: \texttt{figures/tile\_audit\_cost.png}.\end{minipage}}
\caption{Stylized monitor bandwidth: tile-backed logs amortize auditing as the number of checked entries grows.}
\label{fig:tile_cost}
\end{figure}


\paragraph{Inflation under retries.}
Figure~\ref{fig:stop_when_pass} shows stop-when-pass inflation.
Figure~\ref{fig:sweep_grinding} shows that sweeping across many evaluation definitions inflates false approvals toward $1$.

\paragraph{Commitment timing attacks.}
Figure~\ref{fig:commit_timing} shows that post-hoc commitment (choosing the beacon round after seeing outcomes) defeats the intent of public randomness.

\begin{figure}[t]
\centering
\fbox{\begin{minipage}{0.92\linewidth}\centering\small Figure omitted in source-only release. Requested file: \texttt{figures/commit\_timing\_attack\_beacon.png}.\end{minipage}}
\caption{Commitment timing attacks: choosing the beacon round after partial observation enables bias.}
\label{fig:commit_timing}
\end{figure}

\paragraph{Grinding and partial leakage.}
Figure~\ref{fig:abort} and Figure~\ref{fig:partial_leak} show that aborts and partial leakage must be treated as spend.

\begin{figure}[t]
\centering
\fbox{\begin{minipage}{0.92\linewidth}\centering\small Figure omitted in source-only release. Requested file: \texttt{figures/abort\_partial\_leakage.png}.\end{minipage}}
\caption{Partial leakage (seeing a noisy proxy of the outcome) can also enable grinding; log-and-spend blocks it.}
\label{fig:partial_leak}
\end{figure}

\appendix
\section{Checkpoint and proof material (implementation sketch)}
To keep the paper tight, we only sketch the log artifacts needed for independent verification.
A practical implementation can follow tile-backed transparency-log specifications and libraries (e.g., \cite{c2sp_tlog_checkpoint,c2sp_tlog_tiles,c2sp_tlog_proof,c2sp_tlog_witness,tessera_github}).

\paragraph{Checkpoint.}
A checkpoint is a signed statement of the log's current tree state (log identifier, tree size, and root hash), optionally with witness cosignatures.
In tile-backed ecosystems, checkpoints are designed to be \emph{cheap to fetch} and \emph{cheap to compare} across clients and witnesses.

\paragraph{Smallest auditable checkpoint tuple.}
For the maintained worked bundle, the smallest non-decorative tuple worth stating in prose is:
\begin{quote}\small
\nolinkurl{log_id=example-transparency-log, checkpoint=tree_size=42 root=EXAMPLE_ONLY}\\
\texttt{proof\_pointer=consistency://example-transparency-log/41-42}\\
\texttt{receipt\_digest=095e828fd09ef186757f5dd49e84769603a2b59ea77c4a8bfb0edb720bc64cc1}\\
\texttt{bundle\_id=bundle://worked-example/fallback-cct, support\_pointer=support://fallback/notarized-attempt-log.example.}
\texttt{release\_id=worked-example-draft-199, artifact\_note\_version=1.99, plan\_spec=sha256:e9168dbec2cabd9115a7f3ea9aebe0a75ff910d2295e450a340716cc7e541193.}
\end{quote}
A later terse paper may quote that tuple only when the surrounding public wrapper changed but the support packet itself did not.
If one of those fields changes, the support packet has changed and the old tuple should fail closed rather than masquerading as the same checkpoint-hardening note.
The local witness-policy template below is an optional sidecar for operator assumptions, not part of this smallest shipped worked-route tuple.

\paragraph{Inclusion proof.}
For a manifest entry, the publisher supplies either:
(i) a compact inclusion proof against a referenced checkpoint, or
(ii) the tile path sufficient for a monitor to reconstruct the inclusion proof from cached tiles.
In practice, monitors fetch the latest checkpoint and the necessary tiles; proofs are then recomputable offline.

\paragraph{Witnessing and cross-logging.}
Witness protocols allow third parties to cosign checkpoints after verifying append-only consistency, while cross-logging submits checkpoints to an independent log for additional audit surface. Public witness directories can simplify bootstrap and diversity (e.g., via public witness lists in the transparency-logs ecosystem) \cite{c2sp_tlog_witness}.
Both mechanisms are compatible with keeping the underlying artifacts off-chain: only hashes and manifests are logged.


\section{Witness policy template (minimal)}
To keep artifacts compact, you can publish a \emph{witness policy} alongside the log identifier and checkpoint origin.
Below is a minimal template (informal) that is sufficient for third parties to reproduce acceptance conditions.

\begin{verbatim}
log_id: "anon-dht-eval-log"
checkpoint_origin: "anon-dht-eval/v1"
recognized_witnesses:
  - name: "w1"  pubkey: "..."
  - name: "w2"  pubkey: "..."
  - name: "w3"  pubkey: "..."
  # ... total W witnesses
quorum_q: 5
checkpoint_period: "1h"
client_rule:
  accept_checkpoint_if:
    - ">= q witness cosignatures"
    - "checkpoint in append-only log"
  restart_if:
    - "witness set changes"
    - "epoch/beacon id changes"
notes:
  - "Cross-log checkpoints weekly into an independent log"
  - "Distribute witness keys/endpoints via TUF; rotate quarterly"
\end{verbatim}

The policy can be distributed via TUF (or your existing release channel), while the transparency log binds the \emph{manifest hashes}.
This keeps the on-net paper small while still enabling independent verification.
In the stop map above, that witness policy remains only a packet component: it supports an Eval~1 claim, but it does not by itself settle retry-prefix validity, runtime conformance, or public lineage.

\begin{thebibliography}{99}

\bibitem{series:evalnotary}
Anonymous.
\newblock Notarized Anonymity Certificates Under Retries: Grinding-Resistant Evaluation via Transparency Logs, Public Randomness, and $\delta$-Spending.
\newblock Evaluation series Paper~1, 2026. (This archive.)

\bibitem{series:transparencyprimitives}
Anonymous.
\newblock Transparency Primitives Interface: Beacons, VRFs, and Append-Only Logs.
\newblock Series synthesis note (Synthesis~9), 2026. (This archive.)

\bibitem{series:workedexample}
Anonymous.
\newblock Worked Example: Instantiating a Tiered Reader-Privacy Receipt.
\newblock Synthesis~17, The-One-True-Archive (2026).

\bibitem{series:pairedterminalcutoverrules}
Anonymous.
\newblock Paired Terminal Cutover Rules for Stable Review Spines: When Later Terse Papers Must Stop, Widen, or Reopen.
\newblock Series synthesis note (Synthesis~75), 2026. (This archive.)

\bibitem{series:artifactpolicy}
Anonymous.
\newblock Artifact policy and supporting materials for the anonymity/AnonDHT series.
\newblock 2026.

\bibitem{series:opB}
Anonymous.
\newblock Auditable Runtime Conformance for Anonymous Traces.
\newblock Operational series Paper~B, 2026. (This archive.)

\bibitem{series:abomoinl}
Anonymous.
\newblock ABOM/OINL: Public Claim Objects and Delta Surfaces.
\newblock Anonymity series Paper~C, 2026. (This archive.)

\bibitem{series:releaseA}
Anonymous.
\newblock Release Property Ledgers, Receipts, and Public Lineage.
\newblock Release and Destination series Paper~A, 2026. (This archive.)

\bibitem{c2sp_tlog_checkpoint}
\newblock Transparency Log Checkpoint Format (tlog-checkpoint).
\newblock C2SP specification, 2026.
\newblock Accessed 2026-02
\newblock \url{https://c2sp.org/tlog-checkpoint}

\bibitem{c2sp_tlog_proof}
\newblock Transparency Log Offline Proof Format (tlog-proof).
\newblock C2SP specification, 2026.
\newblock Accessed 2026-02
\newblock \url{https://c2sp.org/tlog-proof}

\bibitem{c2sp_tlog_tiles}
\newblock Transparency Log Tiles API (tlog-tiles).
\newblock C2SP specification, 2026.
\newblock Accessed 2026-02
\newblock \url{https://c2sp.org/tlog-tiles}

\bibitem{c2sp_tlog_witness}
\newblock Transparency Log Witness Protocol (tlog-witness).
\newblock C2SP specification, 2026.
\newblock Accessed 2026-02
\newblock \url{https://c2sp.org/tlog-witness}

\bibitem{cloudflare_azul_ct_2025}
Chris Hartwig and Vincent Van Doordrecht.
\newblock A next-generation Certificate Transparency log built on Cloudflare Workers.
\newblock Cloudflare blog, 2025.
\newblock \url{https://blog.cloudflare.com/azul-certificate-transparency-log/}.

\bibitem{dahlberg2018ctgossip}
Rasmus Dahlberg et al.
\newblock Aggregation-Based Certificate Transparency Gossip.
\newblock \texttt{arXiv:1806.08817}, 2018.

\bibitem{dirksen2021logpicker}
Alexandra Dirksen et al.
\newblock LogPicker: Strengthening Certificate Transparency Against Covert Adversaries.
\newblock \emph{PoPETs}, 2021.
\newblock \url{https://petsymposium.org/popets/2021/popets-2021-0066.pdf}.

\bibitem{hicks2023sok_transparency}
Alexander Hicks.
\newblock SoK: Log Based Transparency Enhancing Technologies.
\newblock \texttt{arXiv:2305.01378}, 2023.

\bibitem{slsa_provenance_v1}
\newblock SLSA Provenance v1.0.
\newblock slsa.dev specification, 2026.
\newblock Accessed 2026-02
\newblock \url{https://slsa.dev/spec/v1.0/provenance}

\bibitem{static_ct_api_spec}
\newblock Static Certificate Transparency API.
\newblock C2SP specification, 2026.
\newblock Spec for tile/static CT endpoints (also mirrored at c2sp.org); accessed 2026-02
\newblock \url{https://c2sp.org/static-ct-api}

\bibitem{syta2016cosi}
Ewa Syta et al.
\newblock Keeping Authorities ``Honest or Bust'' with Decentralized Witness Cosigning.
\newblock In \emph{IEEE S\&P}, 2016.

\bibitem{tessera_github}
\newblock Tessera: A Go library for tile-based transparency logs.
\newblock transparency-dev GitHub repository, 2026.
\newblock Accessed 2026-02
\newblock \url{https://github.com/transparency-dev/tessera}

\bibitem{tesseract_blog_2025}
\newblock Introducing TesseraCT.
\newblock blog.transparency.dev, 2025.
\newblock Announcement of TesseraCT (Static CT API implementation) built atop tile-backed transparency logs; accessed 2026-02
\newblock \url{https://blog.transparency.dev/introducing-tesseract}

\end{thebibliography}

\end{document}
